\subsection{Spread-ray estimate for the two-mark calculation}\label{c4:spread}

The radial series is complex, and the charge form is not positive on
the whole doubled space.  Its convergence therefore requires a bound for
the actual nonnegative screening counts.  We prove a strict quadratic
bound for those counts after adding a controlled part of the positive
slope heat.  The remaining heat will make the transfer links
Hilbert--Schmidt. To specify the data used in the check, take the integral-span basis symbols \(\alpha_{ij}=C^j\alpha_{i0}\) (\(i=1,2;\ j\bmod8\)); their two diagonal Gram blocks have first row \((2,-1,1,0,0,0,1,-1)\), their cross block \((0,0,0,-2,2,-2,0,0)\), all circulant (pass to the nondegenerate quotient). Then \(A=S_{\alpha_{10}}S_{\alpha_{20}}C\). The rays used for the expanded unipotent can be grouped in each sector in the following list \(b_{is}\), in order \(s=1,\dots,5\), writing \(i'\) for the other sector:
\[
 \alpha_{i0}-\alpha_{i2}-\alpha_{i3}-2\alpha_{i'0},\
 \alpha_{i0},\ -\alpha_{i2},\ \alpha_{i7},\
 \alpha_{i7}-\alpha_{i1}-\alpha_{i2}.
\]
These are \(1F,12,24,E3,E4\) in matrix labels. We use angles \(\phi_{is}\) which, up to common translation, are arbitrarily small fixed perturbations of \((-\pi,0,0,0,0)\) in both sectors. Take them distinct, in the required factor order within groups. The invariant part of \(A\) is a simple positive line with orthogonal projection \(\Pi\); the other eigenvalues have absolute value 1.

Write \(H_{ij}(u)\) temporarily for the real kernel of \eqref{c4:Rad-2}, using this \(R=A\) and root types \(i,j\) (flattening the two indices), including angle difference; \(u\) is real-position difference. It is even in \(u\) and matrix-symmetric: in the real part both frequency signs occur with half weight, radial damping uses \(|\Lambda_l|\) (sign by branch), and the angle exponential uses \(-i\Lambda_l(\phi_j-\phi_i)\), using the projections and adjunction as in \eqref{c4:Rad-2}. The transform (frequency \(\xi\), no \(2\pi\) in the exponential) satisfies
\[
 \widehat H_{ij}(\xi)+V_{ij}/\xi^2
 =\frac{\pi}{|\xi|} M_{ij}(\xi),\qquad V_{ij}=(b_i,\Pi b_j),\qquad
 M_{ij}=\sum_{k\in\mathbb Z}(b_i,A^k b_j)e^{-|\xi||\phi_j-\phi_i-2\pi k|}.
\]
In fact the exponential sum on the right is the periodized transform for the shifted frequencies \(\gamma+\mathbb Z\), since \(e^{-|\xi||w|}\) has transform \(2|\xi|/(l^2+\xi^2)\); pair \(A^{-k}b_i\) with \(b_j\). As usual analytic calculus near each branch also proves the identity on Jordan blocks. The missing zero frequency gives the invariant correction.

\paragraph{The matrix bound at the unperturbed rays.} The following strict matrix inequality is the main algebraic input. First take unperturbed angles, \(p=e^{-\pi|\xi|}\), and put
\[
 B_0=2\begin{pmatrix}
 0&0&0&0&0\\0&0&0&0&0\\0&0&0&1&1\\0&0&1&0&1\\0&0&1&1&2
 \end{pmatrix},\quad B_1=\tfrac32(2,0,1,1,2)^t(2,0,1,1,2),
\]
and \(N_m\) the two-sector symmetric matrix with diagonal blocks zero and cross block \(B_m\). Then
\begin{equation}
 M-(1-p)(N_0+pN_1)\ \ge\ (1-p) I/20.                                      \label{c4:Spread-1}
\end{equation}
Here are details of one polynomial verification (rational entries throughout). Decompose by sector parity \(\delta=\pm1\). Multiply the difference after subtracting also \((1-p)I/20\) on the right, on that block, by
\[
 D_+=(1-p^2)(p^4+p^2+1)^2(p^4-p^2+1),\qquad
 D_-=(1+p^2)(p^4-p^2+1)^3/(1-p)
\]
respectively. The resulting matrix \(P\) is polynomial of degree \(d=16,15\), respectively. All coefficient matrices \(T_j=[x^j](1+x)^d P(1/(1+x))/\binom dj\) are positive definite, except \(j=0,1\) in the plus block where the unnormalized coefficients are \(16,144\) times the Gram \((1,-1,1,1,2)^t(1,-1,1,1,2)\). For convenience the table gives lower bounds (rounded down to simple rationals) on the leading minors of \(T_j\) excluding just those two:
\[
 \begin{array}{c|ccccc}
  &1&2&3&4&5\\\hline
 +&7/10&1&7/10&1/4&1/50\\
 -&1&3&5&6&4
 \end{array}
\]
To verify entries of the calculation, the following short coefficient formulas suffice. Use \(g_k=((b_{1s},A^k(b_{1t}+\delta b_{2t})))_{s,t=1}^5,\ h_{st}=|{\bf1}_{s=1}-{\bf1}_{t=1}|\), where for this formula order both groups with the first at angle 0 and others at \(\pi\). Form polynomial coefficients \(P_l\) (up to \(d\)) by multiplying \(D_\delta\) in Taylor coefficients at zero against
\[
 \left(\sum_{k\le0}(g_k)_{st}p^{-2k+h_{st}}
       +\sum_{k>0}(g_k)_{st}p^{2k-h_{st}}\right)_{st}
       -(1-p)(\delta B_0+p\delta B_1+I/20).
\]
In this displayed recipe for unequal groups take \(s=1\) first and reflect across the diagonal for the transposed entry. Then the columns of the table are bounds on
\[
 \frac{1}{\binom dj^r}\sum_{\sigma\in S_r}\operatorname{sgn}(\sigma)
 \prod_{a=1}^r\sum_{l=0}^{d-j}\binom{d-l}{j}(P_l)_{a,\sigma(a)}.
\]
Degrees (or cancellation of higher coefficients) follow by summing the two geometric matrix series using \(g_k\) (equivalently \((A^6-I)^3=0\) and sum the first six residues with their first two differences of step 6). This proves \eqref{c4:Spread-1} by positive polynomial evaluation.

\paragraph{Distinct rays and a screened invariant mode.} Let angles now be perturbed by \(\Delta_i\). For sufficiently small fixed \(\omega>0,\kappa>0,|\Delta_i|\), the bound persists, in the form
\begin{equation}
 \widehat H(\xi)+\frac{1-\omega}{\xi^2+\kappa^2} V\ \ge\
 \frac{\pi(1-p)}{|\xi|}\left((N_0+pN_1)\odot(e^{-|\xi||\Delta_i-\Delta_j|})_{ij}+ I/100\right).        \label{c4:Spread-2}
\end{equation}
Indeed at high frequency in units of \(\pi/|\xi|\), only the same-unperturbed-group direct entries survive, multiplied entrywise by the indicated exponential Gram; \(N_0\) also only uses same-group entries. The Gram Schur product retains the limiting strict gap since the Gram diagonal is 1. At low frequency \(\widehat H\) is continuous jointly in angles and frequency by the decaying radial series including the integrable translated-angle log singularities (or the omitted-zero resolvent formula). On \(\ker V\) the limiting gap of \eqref{c4:Spread-1} is strict, while the added large positive matrix dominates on its complement. Compact remaining frequencies follow by continuity.

\paragraph{Coercivity for point configurations.} Consequently on the line, for nonnegative weighted finite root configurations, position ordered with weights \(t_x\), unit cells \(I\), there are constants \(c,C>0\) such that
\begin{equation}
 \sum_{x<y}t_xt_y H_{xy}(u_y-u_x)
 +\frac{1-\omega}{2}\|Y\|_\kappa^2
 \ \ge\ c\sum_I(\sum_{x\in I}t_x)^2-C\sum_x t_x^2,\qquad Y=\sum_x t_x\Pi b_x\delta_{u_x},
                                                                            \label{c4:Spread-3}
\end{equation}
where the squared norm uses the Euclidean metric induced on the invariant line and covariance \((-\partial_u^2+\kappa^2)^{-1}\). The matrices \(N_m\) have nonnegative entries, and the scalar multipliers accompanying those entries have positive inverse kernels. Indeed, for \(0\le a<b\),
\[
 \frac{e^{-a|\xi|}-e^{-b|\xi|}}{|\xi|}
       =\int_a^b e^{-t|\xi|}\,dt,
\]
whose inverse Fourier transform is an integral of positive Cauchy kernels. This applies to \(p^m(1-p)/|\xi|\), also after multiplication by the extra exponential in \eqref{c4:Spread-2}. Since the point weights are nonnegative, these terms can be discarded; the \(I/100\) term retains squared-cell control. No positive-semidefiniteness of the full matrices \(N_m\) is being asserted. To remove diagonal singularities first smooth each point by a Cauchy kernel of small width and apply the bound by Fourier quadratic forms. The smoothed Cauchy inverse kernels retain the local positive gap; discard smoothed \(H\)-diagonals at finite linear-in-\(\sum t_x^2\) cost. For fixed distinct angles (width \(<2\pi\)), each total kernel on the left before smoothing decays exponentially; it is Lipschitz piecewise at zero after subtracting, only on identical types, \(2\log(\sqrt{1+u^2}/|u|)\). Indeed differentiate the radial branch sum in \(|u|\) by geometric series as in \eqref{c4:Rad-2}. Cauchy smoothing this logarithm uses \(2\log|(u+i+i\eta)/(u+i\eta)|\) for width \(\eta\), bounded above by the original plus \(O(\eta/(1+u^2))\). On the remainder (also Lipschitz with decaying derivative at infinity) smoothing costs \(o(1)/(1+u^2)\) by splitting the convolution near and away from \(u\). These errors absorb into the gap, proving \eqref{c4:Spread-3}.

\paragraph{Retaining heat in each transfer link.}
The quadratic bound now controls the root counts.  We next explain why
it can be combined with the link construction without spending all of
the slope heat.  Here are the details of adapting the link bounds to \eqref{c4:Spread-3}. Use the same intra-box and neighboring-box short fractions and Fock dampings as before, with \(f,g\) there (tails equal to a small \(\eta>0\)). On a link source-to-target use heat cost \(\int |m|^2 w^2/2\) on the two boxes, \(m\) the invariant-line slope momentum (jump \(\Pi b\)), \(w=f,g\) respectively. Add to the cost
\[
 M'\eta^2\left(|m_L|^2+|m_R|^2-2|m_J|^2\right)
\]
where these are pre-step boundary slopes at left end, right end (computed from target), and joint. This telescopes, and the split heat sums exactly on cycles. \(M'\) can be fixed sufficiently large before choosing \(\eta\). For the roots weighted by \(w\), \(Y\) differs from the derivative of \(w m\) extended by zero by the regular \(w'm\), the weighted (bounded) mark atoms and the two boundary sources of strength \(\eta m_L,\eta m_R\). Thus using \(\omega>0\), their Yukawa cost in \eqref{c4:Spread-3} is at most
\[
 (1-\omega/2)\int |w m|^2/2 + C\int |m|^2/D^2
       + M'\eta^2(|m_L|^2+|m_R|^2)+O(1)
\]
on bounded numbers/sizes of residual marks, \(D\) the large width. The negative cost at joint is absorbed into the spare heat and squared-cell gap since \(w\) there and over one adjoining unit is bounded below uniformly, and \(|m_J|^2\) is controlled by integrated square and counts squared there. Choose \(\eta\) sufficiently small then \(D\) large. The discrepancy between the \(w_xw_y\) short interactions and actual fractions is absorbed exactly as for the previous Hilbert--Schmidt links (exponential decay at large separation and identical-ray singularity favorable). Thus a squared-cell gap and a positive heat fraction remain, enough for Hilbert--Schmidt bounds, bounded-profile Lipschitz estimates and kernel continuity as before. Holomorphic convergence near every real period (for the continuation from small \(Q_L\), without first taking limits) also follows: repeat unweighted periods on the line using \eqref{c4:Spread-3}, on the circle the Yukawa form being bounded by a small-loss heat estimate via the periodic slope derivative, and scale positions including their period slightly complex as in the previous radial estimate. Oscillator interactions perturb by a small squared-cell cost, and angles to marks can be kept distinct.

\paragraph{Finiteness of the angular expansions.}
The same matrix bound also justifies the coefficientwise angular
conjugations. We prove this consequence separately, since its quadratic
form uses the original indefinite charge pairing, whereas the transfer
links above act on a positive Hilbert space.

Let \(m(\phi)\) be the piecewise constant angular momentum path, with
nonnegative root counts \(k_i\) at the fixed rays \(\phi_i\), and with a
fixed finite list of additional bounded jumps. Extend it by
\(m(\phi+2\pi)=A^{-1}m(\phi)\). Put
\[
 Z=m',\qquad
 \overline m=\frac1{2\pi}\int_0^{2\pi}\Pi m(\phi)\,d\phi.
\]
The total invariant charge of \(Z\) is zero by the gluing condition.
Let \(G_\xi^A\) be the Green kernel of
\((-\partial_\phi^2+\xi^2)^{-1}\) with this twisted boundary condition.
Its pairing against root atoms is the exponential periodization
\(M(\xi)/(2|\xi|)\) computed above. Integration by parts, with boundary
terms cancelling by orthogonality of \(A\), gives
\[
 \lim_{\xi\to0}\langle Z,G_\xi^A Z\rangle
   =\int_0^{2\pi}m(\phi)^2\,d\phi-2\pi\overline m^2.
\]
Here the brackets and squares use the original bilinear form. The
subtracted constant mode lies on the positive invariant line.

For the root-only measure \(Z_{\rm root}\), the Fourier identity gives
its regular Green form explicitly:
\[
 \begin{split}
 \lim_{\xi\to0}\left(
   \langle Z_{\rm root},G_\xi^A Z_{\rm root}\rangle
   -\frac{|\sum_i k_i\Pi b_i|^2}{2\pi\xi^2}\right)
 &=\frac1{2\pi}\sum_{i,j}k_i k_j\widehat H_{ij}(0)\\
 &\ge c\sum_i k_i^2-C\left|\sum_i k_i\Pi b_i\right|^2.
 \end{split}
\]
The inequality is \eqref{c4:Spread-2} at zero frequency; its
entrywise nonnegative terms can be discarded on the nonnegative count
cone. Neutrality bounds \(\sum_i k_i\Pi b_i\) by the additional bounded
jumps. The regular mixed terms with those jumps cost at most
\(C(1+\sum_i k_i)\), while the invariant poles cancel in the total
measure. Absorbing this linear error and restoring the invariant mean
therefore yields
\[
 \int_0^{2\pi}m(\phi)^2\,d\phi
       \ge c\sum_i k_i^2+2\pi\overline m^2-C.
\]
The gluing equation determines the off-invariant entrance momentum as a
fixed linear function of the jumps. Its invariant part differs from
\(\overline m\) by \(O(1+\sum_i k_i)\). After reducing the constants, the
last display thus controls both the squared counts and the squared
entrance momentum in any fixed Euclidean norm on the charge space.

The argument also works with coincident group angles for cut finiteness,
provided all first types are grouped before all other types; the
additional bounded insertions may have any order and angles. For fixed
field powers and cut energy \(H\), all successive energies are bounded
above. The angular average bound consequently gives finite support. For
evaluation at small \(Q_L\), damping by that same average dominates the
exponential-in-squared-count oscillator bounds exactly as in
\eqref{c4:Ang-2}.

During coefficientwise conjugations cutting instead at \(C\), one only needs that the orthogonal-off-invariant \((1-C)^{-1}\) images of the jumps in \(g^{\pm1}\), on their nonnegative cone, control squared counts. Similarly cutting at \(B=S_{\alpha_{10}}C\) during the partial conjugation one uses \((1-B)^{-1}\) images of the neutral roots \(E4,24\) of sector 1 together with \(g_2^{\pm1}\). The comparisons below are lower bounds on the nonnegative count cone: replacing positive off-diagonal entries by zero decreases the quadratic form there, but need not do so on arbitrary vectors. For \(C\) the lists per sector are \(-\alpha_{i0},-\alpha_{i1},-\alpha_{i1}-\alpha_{i2},\alpha_{i7}\). The diagonal of the stated Gram times 16 is \(11,11,16,11\) twice, with within-sector negative entries only at \(12,24\) equal to \(-3\) (and transposes), cross-sector negatives at \(14,41\) equal to \(-5\). After the positive off-diagonal entries are removed, the strict diagonal-dominance margins are \((3,5,16,3)\) in each sector, so the true Gram is bounded below by \(3I/16\) on that cone. For \(B\), in the order indicated (using same list for sector 2), four times the Gram after dropping positive off-diagonals is
\[
 \begin{pmatrix}
24&0&-4&0&0&0\\0&8&-4&0&0&0\\-4&-4&9&0&0&-5\\
0&0&0&9&0&-5\\0&0&0&0&28&-6\\0&0&-5&-5&-6&9
 \end{pmatrix}>0.
\]
To verify positivity, eliminate the first, second, fourth and fifth coordinates. The remaining Schur complement is
\[
 \begin{pmatrix}19/3&-5\\-5&311/63\end{pmatrix},
 \qquad \det=1184/189>0.
\]
These computations use just the circulant and reflections above. Thus squared-cut-momentum bounds hold with any bounded additional shifts; on the \(B\)-side for evaluated traces near small \(Q_L\) one can bunch the neutral-root contours closely.

\paragraph{The suppression profiles.}
We finish by identifying which root activities the large fused charge
suppresses.  The two sector sizes enter only through these profiles;
the coercivity constants above can be chosen uniformly. \((A+A^{-1})h_i=h_i\); for \(b_{js}\), \((x,y)=((h_i,b_{js}),(Ah_i,b_{js}))\) is zero across sectors and within equal respectively to 
\[(1,0),(1,0),(0,0),(0,1),(0,0).\]
As in the single-mark case a charge \(a_i h_i+d_i Ah_i\) at angle and position zero gives real log-activity on its own axis (before images)
\[
 (a_i x+d_i y) F_X+((a_i+d_i)y-d_i x) F_Y,\quad
 F_X=J(\psi)+J(\psi-\pi/3),\quad F_Y=J(\psi+2\pi/3)+J(\psi+\pi/3),
\]
where \(J(\theta)=\log|1-r e^{i\theta}|-\log|1+r e^{i\theta}|,\ r=e^{-|u|/6},\ \psi=\phi_b/6\); this follows by the six offsets in \eqref{c4:Rad-2}, in either position order. For \(d_i/a_i\) close in both sectors to a common \(\rho\in[0,1]\), choose the two group angles near \(-\pi\) and 0. For \(0<\rho<1\), precisely at \(-\pi\), \(F_X=F_Y<0\); precisely at 0, \(F_X<0,F_Y=0\); the inequalities for the required nonneutral types thus have uniform strict sign after division by \(r a_i\). Perturb away slightly, since a singularity at \(r=1\) only improves the cost. For \(\rho=0\) instead shift both angles first very slightly negative (\(F_Y/r\) has strictly positive \(\psi\)-derivative at 0). For \(\rho=1\) shift slightly positive (\((F_X-F_Y)/r\) has strictly negative derivative at \(-\pi/6\)). These statements follow directly from \(J=-\operatorname{atanh}(2r\cos\theta/(1+r^2))\). By neighborhoods and a finite cover one can thus choose distinct angles as needed above with a uniform strict bound \(-c a_i e^{-|u|/6}\), leaving angle 0 free; complex logs decay also in absolute value at the two infinities at that rate scale. In placing the compensator near 0, only the first root type needs to precede it; the later group can be put before or after the bounded-angle-near-zero cluster, using commutation through all external fields as necessary (their pairings there are nonnegative). This explains simultaneous availability of the damping profiles and the spread estimate.
